\section{Discussion}
\label{sec:discussion}

\para{Where does the gradient-masking picture fit?} The hypothesis we started from is that safety tuning learns a defence keyed to surface features. For \emph{vocabulary}, the evidence mostly speaks against it. Refusal follows the meaning of a request in every model we tested, from 0.5B to 32.8B, and removing trigger words costs a moderate 12--16\pp of refusal that mostly turns into sanitised, not harmful, compliance. For \emph{framing}, the picture fits well. The same model that refuses 95.5\% of direct requests complies with 86\% of them when told it is a persona without rules. Within a framing, models track semantics; what they fail to do is recognise the same harmful request once it is wrapped in a fictional or persona frame. This looks like refusal behaviour that is tied to the distribution of the safety training data, which is the essence of masking, but the relevant feature is the conversational frame and not the words. This view is consistent with mismatched generalisation \cite{wei2023jailbroken}: the frames we use are common in pre-training data, but likely rare in safety data.

\para{Why does scale not help?} Our scale ladder spans six families, so family and size are confounded. The cleanest evidence is the Qwen2.5 ladder, where fiction compliance rises from 30\% at 0.5B to 76\% at 3.1B, drops to 11\% at 14.8B, and rises again to 40\% at 32.8B. Across families, Llama-2-7B-chat and Llama-3.2-1B/3B resist every framing, while Llama-3.1-8B and Gemma-2-9B break under the persona prompt. A plausible reading is that robustness to framing depends on whether the safety data of a given checkpoint happened to cover that frame. Greater capability can even raise risk, as in \etwo, where large models comply more in low-resource languages because they understand the request.

\para{Implications for red-teaming and evaluation.} Three practical lessons follow. First, red-teaming of open models should prioritise cheap framing attacks such as fiction and personas, because a single fixed template already exposes most of the blind spots we found, and the vulnerable frame differs by family. Second, multilingual jailbreak rates for small models should be reported as harmful compliance, not as non-refusal; otherwise they are inflated three- to ten-fold by incoherent replies. Third, lexical robustness should be measured with a keyword-kept control and with a harm check on the rewrite, since in our data the apparent effect shrinks by about 30\% (from 16.3 to 11.8\pp) once rewrites that sound milder are excluded.

\para{Limitations.} Our labels come from one LLM judge. It was validated on 37 English hand labels and 16 multilingual spot checks, and cross-checked against \harmbench and a string-match detector only in English; non-English refusal recall is imperfect, so non-refusal in \etwo is an upper bound. The paraphrases come from one LLM, and meaning drift is controlled by an LLM equivalence judge and an LLM harm rater, not by humans; the trigger lexicon is fixed a priori and incomplete. We use a single decoding setting (greedy, 256 tokens, no system prompt); system prompts or sampling change refusal rates, and truncation at 256 tokens can hide whether a fictional frame eventually delivers harmful detail. Scale is confounded with family. Models of 7B or more were run on five of nine non-English languages, and each condition was run once, without seed variance. We study open-weight models only, because no API access was available. Finally, \multijail prompts vary widely in severity (some are merely rude), and English refusal rates on them range from 29\% to 75\%.

\para{Broader impact.} This work describes attacks that are already public (fiction framing, personas, translation, and re-wording), and we report aggregate rates without reproducing operational content; all quoted responses are truncated. We believe the main effect of our findings is defensive: they tell model developers which gaps are large and which are not, and they warn that common multilingual metrics overstate risk for small models while understating how risk grows with capability.
